\chapter{Entropy Rate and the Identity $h = \log(1/R)$}
\label{ch09:top}

One of the cleanest bridges between information theory and analytic combinatorics
is the identity
\begin{equation}
  \label{eq:h-log-R}
  h = \log_2(1/R),
\end{equation}
where $h$ is the entropy rate of a stationary source and $R$ is the radius of
convergence of a suitable counting-type generating function. This chapter
develops the identity carefully, warns against a common conflation, and surveys
what is known about entropy rates for natural language and LLM outputs.

%----------------------------------------------------------------------
\section{Entropy Rate}
%----------------------------------------------------------------------

Let $X_1, X_2, \ldots$ be a stationary stochastic process over a finite alphabet
$\Sigma$. The \emph{block entropy} at length $n$ is
\[
  H(X_1^n) := -\sum_{w \in \Sigma^n} p(w)\log_2 p(w),
\]
where $p(w) = \Pr(X_1 \cdots X_n = w)$. The \emph{entropy rate} is
\[
  h := \lim_{n \to \infty} \frac{H(X_1^n)}{n},
\]
when the limit exists. Stationarity guarantees existence and the further
identity
\[
  h = \lim_{n \to \infty} H(X_n \mid X_1^{n-1}).
\]
The conditional form shows that $h$ measures the irreducible uncertainty per
new symbol given arbitrarily long past context.

Informally, $h$ bits per symbol are necessary and sufficient (on average) to
losslessly encode a long sample. For a memoryless uniform source on $\Sigma$,
$h = \log_2|\Sigma|$, the maximum possible. In general,
$0 \le h \le \log_2|\Sigma|$, with $h = \log_2|\Sigma|$ if and only if the
source is uniform i.i.d.

%----------------------------------------------------------------------
\section{Shannon--McMillan--Breiman and Typical Sets}
%----------------------------------------------------------------------

The fundamental limit theorem of source coding is the following.

\begin{theorem}[Shannon--McMillan--Breiman]
  Let $\{X_n\}$ be an ergodic stationary source over $\Sigma$ with entropy
  rate $h$. Then
  \[
    -\frac{1}{n}\log_2 p(X_1 \cdots X_n) \;\to\; h
    \quad \text{almost surely as } n \to \infty.
  \]
\end{theorem}

This is the \emph{asymptotic equipartition property} (AEP). The block
log-probability concentrates, almost surely, at $-nh$.

\begin{definition}[Typical set]
  For $\epsilon > 0$, the \emph{$\epsilon$-typical set} at length $n$ is
  \[
    T_n^{(\epsilon)} := \Bigl\{ w \in \Sigma^n :
      \bigl| {-\tfrac{1}{n}\log_2 p(w)} - h \bigr| < \epsilon \Bigr\}.
  \]
\end{definition}

The AEP immediately yields three properties: (i) $\Pr(X_1^n \in T_n^{(\epsilon)})
\to 1$; (ii) each $w \in T_n^{(\epsilon)}$ satisfies
$2^{-n(h+\epsilon)} \le p(w) \le 2^{-n(h-\epsilon)}$; and therefore
(iii) the cardinality satisfies
\[
  (1-\delta_n) \cdot 2^{n(h-\epsilon)}
  \;\le\; |T_n^{(\epsilon)}|
  \;\le\; 2^{n(h+\epsilon)}
\]
for any $\delta_n \to 0$. Setting $\epsilon \to 0$ slowly: $|T_n| \approx 2^{nh}$
up to subexponential factors.

%----------------------------------------------------------------------
\section{From Typical Sets to Generating Functions}
%----------------------------------------------------------------------

Package the cardinality of the typical set into an ordinary generating function.
Let $a_n := |T_n|$ for a fixed (small) $\epsilon$. Consider
\[
  A(z) = \sum_{n \ge 0} a_n z^n, \qquad a_n \approx 2^{nh}.
\]
By the Cauchy--Hadamard formula (Chapter 2), the radius of convergence is
\[
  \frac{1}{R} = \limsup_{n\to\infty} a_n^{1/n} = 2^h.
\]
Therefore $R = 2^{-h}$ and we recover the identity \eqref{eq:h-log-R}:
\[
  \boxed{h = \log_2(1/R).}
\]

\begin{remark}
  The identity is classical for finite-state sources. If the source is generated
  by a finite automaton with transition matrix $M$, then $a_n$ counts paths of
  length $n$ in the automaton. The transfer-matrix method (Chapter 7) gives
  $[z^n]A(z) \sim C \cdot \rho(M)^n$, where $\rho(M)$ is the Perron eigenvalue.
  The radius of convergence is $R = 1/\rho(M)$, and the entropy rate of the
  corresponding stationary Markov chain satisfies $h = \log_2 \rho(M)$.
  The identity thus unifies the Perron--Frobenius spectral picture (Chapters 6--7)
  with the information-theoretic picture.
\end{remark}

%----------------------------------------------------------------------
\section{A Crucial Distinction}
%----------------------------------------------------------------------

The identity $h = \log_2(1/R)$ applies to a \emph{counting} generating function
whose coefficients grow like $2^{nh}$. It does not apply to the
\emph{length probability} generating function, and confusing the two leads to
error.

Define the length probability generating function by
\[
  P(z) := \sum_{n \ge 0} p_n\, z^n, \qquad p_n = \Pr(|W| = n),
\]
for a model that generates strings of random length (such as an autoregressive
language model with an end-of-sequence token). Since the $p_n$ are probabilities
summing to 1, we have $P(1) = 1$, and therefore the radius of convergence satisfies
$R_P \ge 1$. Consequently,
\[
  \log_2(1/R_P) \le 0,
\]
which cannot encode a positive entropy rate. What $R_P$ and the singularity
structure of $P$ near $z = 1$ \emph{do} encode is the tail behavior of the
length distribution: exponential decay, subexponential decay, polynomial decay,
and so on.

Two distinct generating functions thus play two distinct roles.
\begin{enumerate}
  \item \textbf{The counting-type generating function} $A(z) = \sum a_n z^n$,
    where $a_n$ counts strings of length $n$ in the typical set (or, for
    finite-state sources, paths of length $n$ in the underlying automaton).
    Its radius satisfies $R = 2^{-h}$ and encodes the entropy rate.
  \item \textbf{The length probability generating function} $P(z) = \sum p_n z^n$.
    Its radius satisfies $R_P \ge 1$ and its singularity type encodes the
    length-distribution tail.
\end{enumerate}

When you encounter the phrase ``the generating function of the LLM'' in a
paper, you must ask which of these two objects is intended. The two questions
one can ask about a source --- how large is the typical set, and how are string
lengths distributed --- correspond to two generating functions with radii of
convergence on opposite sides of $z = 1$.

%----------------------------------------------------------------------
\section{Classical Entropy-Rate Estimation for Natural Language}
%----------------------------------------------------------------------

Entropy-rate estimation for natural language has a history of more than seventy
years, and the problem remains open at the precision needed for modern
applications.

\paragraph{Shannon (1951).}
Shannon \cite{Shannon1951} estimated the entropy of printed English via a
guessing experiment. Human subjects predicted the next character of a text
given the preceding characters; the distribution of guesses yields an upper
bound on conditional entropy, while the reciprocal of the guessed character's
rank in a natural ordering yields a lower bound. For contexts of length
$N = 100$, Shannon's estimates placed the entropy rate in the range
$[0.6, 1.3]$ bits per character. In the language of Section~3, Shannon was
estimating $\log_2(1/R)$ for the counting generating function of English,
though he never wrote that function down. The width of his interval reflects
both experimental variance and the slow convergence of the conditional entropy
to its limit, a theme we return to below.

\paragraph{Kontoyiannis et al.\ (1998).}
Kontoyiannis, Algoet, Suhov, and Wyner \cite{KontoyiannisEtAl1998} sharpen
the estimation problem using longest-match-length statistics. Let $L_n$ denote
the length of the longest prefix of $(X_{n+1}, X_{n+2}, \ldots)$ that has
appeared as a substring in $(X_1, \ldots, X_n)$. Their estimator is
\[
  \hat{h}_n := \frac{\log_2 n}{L_n}.
\]
Under a Doeblin-type mixing condition on the source, $\hat{h}_n \to h$ almost
surely. The convergence rate is controlled by how cleanly the dominant
singularity of $A(z)$ is separated from the remainder of the spectrum on the
circle $|z| = R$; a simple isolated pole gives geometric convergence, while
a more complex singularity structure degrades the rate.

\paragraph{Szpankowski (2001).}
The book \cite{Szpankowski2001} applies the full Flajolet--Sedgewick toolkit
--- Mellin transforms, depoissonization, singularity analysis --- to these
questions in the context of digital data structures and source models. Part V,
Chapter 17 of the present text covers this material in depth.

%----------------------------------------------------------------------
\section{Modern LLM-Era Estimates}
%----------------------------------------------------------------------

\paragraph{Takahashi and Tanaka-Ishii (2018).}
Takahashi and Tanaka-Ishii \cite{TakahashiTanakaIshii2018} extrapolate neural
language-model cross-entropies simultaneously to infinite training data and
infinite context length. Their extrapolation procedure fits the observed
cross-entropies as a function of both corpus size and context window using a
two-parameter model, and takes the joint limit. For English, they obtain an
entropy-rate estimate of approximately 1.12 bits per character. In the present
framework, this is an estimate of $\log_2(1/R)$ for the English counting
generating function, computed indirectly via the cross-entropy of a neural
model serving as a proxy for the true source distribution.

The result is consistent with, and more precise than, Shannon's original
interval. The narrowing was achieved not by a more careful experiment but by
using a much better predictor, trained on vastly more text with much longer
effective context.

\paragraph{Scheibner, Smith, and Bialek (2025).}
Scheibner, Smith, and Bialek \cite{ScheibnerSmithBialek2025} return to
Shannon's guessing-game paradigm but replace human subjects with frontier large
language models as predictors. The central finding is that conditional entropy
continues to decrease with context length out to $N \sim 10^4$ characters, at
scales of dependency that Shannon could not access and that substantially
exceed the context lengths used in \cite{TakahashiTanakaIshii2018}. The
running entropy-rate estimate does not settle within the range of context
lengths they study.

%----------------------------------------------------------------------
\section{Slow Convergence as a Diagnostic for Singularity Type}
%----------------------------------------------------------------------

The identity $h = \log_2(1/R)$ specifies the location of the dominant
singularity of $A(z)$ but says nothing about its \emph{type}. The type of the
singularity determines the rate at which the conditional entropy
$H(X_{n+1} \mid X_1^n)$ converges to $h$, and this convergence rate is exactly
what entropy-estimation experiments measure.

To be precise, recall from Chapter 4 that the coefficient asymptotics of a GF
with a dominant singularity at $z = R$ depend on the local behavior of $A(z)$
near $R$.

\begin{itemize}
  \item If $A(z)$ has a \emph{simple pole} at $R$, then
    $[z^n]A(z) = C R^{-n}(1 + O(\theta^n))$ for some $\theta < 1/R$.
    The conditional entropy satisfies $H(X_{n+1}\mid X_1^n) - h = O(\theta^n)$:
    exponentially fast convergence. A few hundred characters of context suffice
    to pin down $h$ to three decimal places. This is the Markov chain and
    finite-automaton case.

  \item If $A(z)$ has an \emph{algebraic branch point} of exponent $\alpha$,
    meaning $A(z) \sim C(1-z/R)^\alpha$ as $z \to R$ with $\alpha \notin \mathbb{Z}$,
    then by the Transfer Lemma (Chapter 3),
    \[
      [z^n]A(z) \;\sim\; C' \cdot R^{-n} \cdot n^{-\alpha-1}.
    \]
    The conditional entropy converges with a polynomial residual,
    $H(X_{n+1}\mid X_1^n) - h \sim C'' n^{-\alpha-1}$.

  \item If $A(z)$ has a \emph{logarithmic singularity},
    $A(z) \sim C(1-z/R)^\alpha \log(1-z/R)$, the residual is $(\log n)^{-\beta}$,
    decaying more slowly than any power law.

  \item A \emph{confluence of singularities} on the circle $|z| = R$ introduces
    oscillatory corrections on top of the algebraic or logarithmic envelope.
\end{itemize}

The Shannon-to-Takahashi-Tanaka-Ishii span of $\sim 0.6$ bits is not
measurement noise; it is a signal. Estimates of $h$ via finite-context
predictors converge to the true entropy rate from above as context length
increases, and the \emph{rate} at which they converge encodes the singularity
type of $A(z)$. If the convergence were geometric, the 1951 and 2018 estimates
would agree to many significant figures. They do not. Scheibner et al.'s
finding that the estimate continues to drift downward past $N = 10^4$ characters
is consistent with a singularity that is algebraic or logarithmic, with a slow
polynomial or logarithmic tail on the convergence.

\begin{remark}
  This reframes a classical problem in information theory. Decades of work have
  been devoted to estimating the exponent $h = \log_2(1/R)$. The type of the
  singularity of $A(z)$ --- which controls how fast the estimate converges and
  thus how long a memory the source requires --- has received almost no
  systematic attention. In the language of Part V, this is a question about the
  \emph{nature} of the dominant singularity, not merely its \emph{location}. The
  entropy-rate convergence curves produced by modern LLMs, evaluated at
  progressively longer contexts, are, in principle, a fine-resolution probe of
  that singularity type. Threading together the analytic-combinatorics tools of
  Part I with the estimation procedures of Part V is a program
  taken up in Chapter 17.
\end{remark}

\bigskip

We have completed the foundations of Part II. The identity $h = \log_2(1/R)$
ties together the entropy rate of a stationary source, the asymptotics of its
typical set, and the singularity analysis of its counting generating function.
The distinction between the counting GF and the length probability GF is
essential and is frequently missed. The convergence rate of entropy estimators
to $h$ is a window onto the singularity type of $A(z)$, and that type is what
modern LLM experiments are beginning to resolve.

Part III turns to autoregressive language models as mathematical objects:
first as probability distributions over $\Sigma^*$, then through the lens of
formal-language theory applied to transformer architectures, and finally as
weighted-finite-automaton approximations whose spectral structure directly
instantiates the framework developed in these pages.

\section*{Further Reading}

\begin{itemize}
  \item \cite{Shannon1951} --- the original guessing-game experiment estimating the entropy rate of printed English; the historical starting point for all subsequent entropy-rate estimation work.
  \item \cite{KontoyiannisEtAl1998} --- establishes the longest-match-length estimator $\hat{h}_n = (\log_2 n)/L_n$ and proves its almost-sure convergence to the entropy rate under mixing conditions; the rigorous nonparametric approach to entropy-rate estimation.
  \item \cite{Szpankowski2001} --- applies Mellin transforms, depoissonization, and singularity analysis to entropy-related combinatorial problems; the reference for the analytic toolkit connecting generating-function singularity type to estimation convergence rates.
  \item \cite{TakahashiTanakaIshii2018} --- extrapolates neural language model cross-entropies to infinite data and infinite context, obtaining a precise entropy-rate estimate of approximately $1.12$ bits per character for English.
  \item \cite{ScheibnerSmithBialek2025} --- uses frontier LLMs as predictors in Shannon's guessing paradigm, finding that conditional entropy continues to decrease at context lengths well beyond those studied in earlier work and providing evidence for a slow-converging singularity in the English counting GF.
\end{itemize}
